\documentclass[10pt]{article}

\usepackage[utf8]{inputenc}

\usepackage[T1]{fontenc}

\usepackage{amsmath}

\usepackage{amsfonts}

\usepackage{amssymb}

\usepackage[version=4]{mhchem}

\usepackage{stmaryrd}

\usepackage{bbold}

\usepackage{CJKutf8}

\usepackage{graphicx}

\usepackage[export]{adjustbox}

\graphicspath{ {./images/} }



\begin{document}

как из характеристических, так и нехарактеристических определенным образом ориентированных многообразиӥ.



Для гиперболич. уравнений 2-го порядка носителем начальных данных при постановке смешанной задачи является пространственно ориентированная часть границы $\partial D$. На временны́м образом ориентированной части $\partial D$, как правило, задаются краевые условия такого же типа, как и для параболич. уравнений (см. Смешанная и краевая задачи для параболических уравнений и систем).



Пусть $\Omega$ - область пространства $\mathbb{R} n$ точек $x=\left(x_{1}\right.$, $\left.x_{2}, \ldots, x_{n}\right)$ с достаточно гладкой границей $\partial \Omega$, а



\[

D=\left\{\left(x, x_{0}\right): x \in \Omega, 0<x_{0}<T\right\}

\][



S=\textbackslash left\{\textbackslash left(x, x\_\{0\}\textbackslash right): x \textbackslash in \textbackslash partial \textbackslash Omega, 0<x\_\{0\}<T\textbackslash right\} \textbackslash text \{, \}

]



$\Omega_{0}=\left\{\left(x, x_{0}\right): x \in \Omega, x_{0}=0\right\}, \Omega_{T}=\left\{\left(x, x_{0}\right): x \in \Omega, x_{0}=T\right\}$.



В области $D$ задано линейное гиперболич. уравнение 2-го порядка



\[

\begin{gathered}

u_{x_{0} x_{0}}-a^{i j}\left(x, x_{0}\right) u_{x_{i} x_{j}}+a^{i}\left(x, x_{0}\right) u_{x_{i}}+ \\

+a^{0}\left(x, x_{0}\right) u_{x_{0}}+a\left(x, x_{0}\right) u=f\left(x, x_{0}\right),

\end{gathered}

\]



где подразумевается суммирование от 1 до $n$ по повторяІощимся индексам $i, j$ и форма $a^{i j}\left(x, x_{0}\right) \xi_{i} \xi_{j}$ положительно определена.



Основные смешанные задачи для уравнения (1) охватываются следующей постановкой: в области $D$ найти решение $u=u\left(x, x_{0}\right)$ уравнения (1), удовлетворягощее на $\Omega_{0}$ начальным условиям



\[

\left.u\right|_{x_{0}=0}=\tau(x),\left.\quad u_{x_{0}}\right|_{x_{0}=0}=v(x),

\]



а на $S$ - одному из краевых условий



\[

\begin{gathered}

\left.u\right|_{s}=\varphi_{1}\left(x, x_{0}\right) ; \\

\partial u / \partial N I_{s}=\varphi_{2}\left(x, x_{0}\right) ; \\

\left.\left(\partial u / \partial N+\sigma\left(x, x_{0}\right) u\right)\right|_{s}=\varphi_{3}\left(x, x_{0}\right),

\end{gathered}

\]



где $N$ - конормаль относительного оператора $a^{i j} \frac{\partial}{\partial x_{i}} \frac{\partial}{\partial x_{j}}$. Задачи (2), (3); (2), (4) и (2), (5) принято соответственно называть п е р в ой, в т о р о й и т р е ть е й с ме. III $\mathbf{a}$ н н о й 3 а д а ч е й для уравнения (1).



Для уравнения (1) при довольно общих предположениях относительно өго коэффициентов и границы $\partial D$, а также для заданных функций доказаны существование и единственность как регулярных, так и обобщенных решений всех трех смешанных задач, исследованы структурные и дифференциальные свойства әтих решений в замкнутой области $\bar{D}$ в зависимости от гладкости ее границы [8]. При $n=1$ решение смешанных задач выписывается в явном виде.



Смешанные задачи исследованы для широкого класса линейных и нелинейных гиперболич. уравнений и систем (см. Квазилинейные гиперболические уравления u cucme.nь). Построена удовлетворительная теория смешанных задач для строго гиперболич. уравнений и систем вида



\[

u_{x_{0} x_{0}}=\sum_{|\alpha| \leqslant m} a_{\alpha}\left(x, x_{0}\right) \frac{\left.\partial^{\mid \alpha}\right|_{u}}{\partial x_{1}^{\alpha_{1}} \partial x_{2}^{\alpha_{2}} \ldots \partial x_{n}^{\alpha_{n}}}+f\left(x, x_{0}\right)

\]



с начальными данными на (шространственно ориентированної) части границы области $D$, лежащей на плоскости $x_{0}=0$. Определенный успех лостигнут и при изучении смешанных задач для гиперболич. уравнений и систем в случае, когда носители начальных илии краевых условий представляют собой поверхности вырождения тина или порялка этих уравнениӥ (см. Вырожденное уравнение с частными производными). Наиболее существенные результаты получены для линейных уравнений 2-го порядка вида



\[

\begin{gathered}

u_{x_{0} x_{0}}-\frac{\partial}{\partial x_{j}}\left[a^{i j}\left(x, x_{0}\right) u_{x_{i}}\right]+b^{i}\left(x, x_{0}\right) u_{x_{i}}+ \\

+b^{0}\left(x, x_{0}\right) u_{x_{0}}+c\left(x, x_{0}\right) u=f\left(x, x_{0}\right)

\end{gathered}

\]



с коәффициентами, удовлетворяющими условию



\[

a^{i j} \xi_{i} \xi_{j} \geqslant \lambda_{0} \mu_{0} x_{0}\left(b^{i} \xi_{i}\right)^{2}-\mu_{0} a_{x_{0}}^{i j} \xi_{i} \xi_{j}\left(\forall \xi \in \mathbb{R}^{n}\right),

\]



где $\lambda_{0}$ и $\mu_{0}$ - нек-рые положительные постоянные, и особенно для уравнений вида



\[

y^{m} u_{x x}-u_{y y}+a(x, y) u_{x}+b(x, y) u_{y}+

\]



$+c(x, y) u=f(x, y), x_{1}=x, x_{0}=y \geqslant 0, m=$ const. (6)



\begin{CJK}{UTF8}{mj}下\end{CJK} смешанным задачам с внутренними или внешними краевыми условиями (см. Внешнял $u$ внутрениял краевье задачи) редуцируются математич. модели многих процессов теории рассеяния волн на препятствиях. Напр., к условию излучения Зоммерфельда (см. Излучения условия) приводит задача отыскания решения $u$ волнового уравнения



\[

\square u \equiv u_{x_{0} x_{n}}-u_{x_{1} x_{1}}-\ldots-u_{x_{n} x_{n}}=0

\]



для всех точек $x \in \mathbb{R} n$, лежащих вне ограниченной области $\Omega \subset \mathbb{R} n$, если известно, что производная $u$ по направлению внешней нормали и $\partial \Omega$ обращается в нуль для любого момента времени $x_{0} \geqslant 0$, а начальные условия соответствуют плоской волне, идущей из бесконечности в направлении оси $x_{1}$.



Основными краевыми задачами для гиперболич. уравнений и систем являются Гурса, Дарбу - Пикара и их многомерные аналоги (см. Гурса задача, Коши характеристическая задача, а также [1]).



Задачи Гурса, Дарбу - Пикара и их различные обобщения хоропо исследованы для гиперболич. уравнений и систем 2-го порядка с расщепленными главными частями вида



$u_{x x}-u_{y y}+a(x, y) u_{x}+b(x, y) u_{y}+c(x, y) u=f(x, y)$, где $a, b$ и $c$ - заданные действительные $(m \times m)$-матрицы, $f$ - заданный, а $u$ - искомый $m$-мерные векторы. Существенные результаты получены и для довольно широкого класса гиперболич. систем уравневий 2-го порядка с нерасщепленными главными частями при отсутствии параболич. вырождения. Обнаружен факт неединственности решения характеристич. задачи Гурса $u_{i}(x, x)=\varphi_{i}(x), u_{i}(x,-x)=\psi_{i}(x), x \geqslant 0, i=1,2$, для гиперболич. системы



\[

\frac{\partial^{2} u_{i}}{\partial x^{2}}-\frac{\partial^{2} u_{i}}{\partial y^{2}}+2 \frac{\partial^{2} u_{j}}{\partial x \partial y}=0, i=1,2, i+j=3,

\]



с двумя независимыми переменвыми $x, y$ и найден эффект влияния младших членов на корректность этой задачи [3]. Достаточно полно изучен вопрос влияния характера параболич. вырождений на корректность как локальных, так и нелокальных краевых задач для вырождающихся гиперболич. уравнений и систем [3]. В частности, исследованы основные (локальные) краевые задачи для линейных вырождаюццихя гиперболич. уравнений вида



\[

\begin{gathered}

u_{y y}-k(x, y) u_{x x}+a(x, y) u_{x}+b(x, y) u_{y}+ \\

+c(x, y) u=f(x, y)

\end{gathered}

\]



\includegraphics[max width=\textwidth, center]{2023_07_09_20d5871e1314cef8ae56g-1}

гладкой границеї, установлен факт влияния порядка нехарактеристич. вырождения на корректность задачіІ Дарбу и неравноправия характеристик как носителеї краевых условий [10].



В связи с проблемой поиска многомерных аналогов задач Дарбу и Трикоми (см. Смешанного типа урав-





\end{document}